\documentclass[11pt]{article}
\usepackage[T1]{fontenc}
\usepackage[utf8]{inputenc}
\usepackage{lmodern}
\usepackage[a4paper,margin=27mm]{geometry}
\usepackage{amsmath,amssymb,amsthm,mathtools}
\usepackage{booktabs}
\usepackage{microtype}
\usepackage{xcolor}
\usepackage{enumitem}
\usepackage{listings}
\usepackage[hidelinks]{hyperref}

\definecolor{leanblue}{RGB}{31,78,121}
\definecolor{codegray}{RGB}{245,246,248}
\lstdefinelanguage{Lean}{
  morekeywords={theorem,def,by,have,exact,refine,rw,calc,let,noncomputable,
    namespace,end,import,open,if,then,else,fun,forall,exists},
  sensitive=true,
  morecomment=[l]{--},
  morecomment=[s]{/-}{-/},
  morestring=[b]"
}
\lstset{language=Lean,basicstyle=\ttfamily\small,backgroundcolor=\color{codegray},
  keywordstyle=\color{leanblue}\bfseries,commentstyle=\color{gray},
  columns=fullflexible,breaklines=true,frame=single,rulecolor=\color{gray!30},
  xleftmargin=1mm,xrightmargin=1mm}

\newtheorem{theorem}{Theorem}[section]
\newtheorem{proposition}[theorem]{Proposition}
\newtheorem{corollary}[theorem]{Corollary}
\theoremstyle{definition}
\newtheorem{definition}[theorem]{Definition}
\theoremstyle{remark}
\newtheorem{remark}[theorem]{Remark}

\newcommand{\SU}{\mathrm{SU}(2)}
\newcommand{\Haar}{\mu_{\mathrm H}}
\newcommand{\chiN}[1]{\chi_{#1}}
\newcommand{\K}{K}
\newcommand{\Proj}{\mathsf P}

\title{An Exact $\SU$ Heat-Mode Witness for Transport\\
\large A mechanized bridge from Haar convolution to a non-vacuous finite Dobrushin endpoint}
\author{Lluis Eriksson}
\date{August 4, 2026}

\begin{document}
\maketitle

\begin{abstract}
We give a machine-checked bridge between two independently formalized pieces of
two-dimensional Yang--Mills transport.  The first is the concrete $\SU$ heat
kernel on normalized Haar probability, already constructed as an infinite
character series.  The second is a finite Dobrushin transfer theorem whose
premises had previously been discharged by an exact two-holonomy witness.  The
new result proves that Haar convolution by the infinite heat kernel acts on
every Chebyshev character $\chi_n$ with the exact Casimir eigenvalue
\[
  \exp\!\left(-\frac{t\,n(n+2)}4\right).
\]
It then proves that $\chi_1$ is a nonzero sharp slowest non-vacuum mode, with
rate $e^{-3t/4}$, and identifies this rate definitionally with the parameter
used by the finite witness.  The infinite identity is not postulated: finite
character orthogonality is integrated first and the limit is passed through
the Haar integral by dominated convergence.  The finite carrier consists of
two distinct bundled elements of $\SU$, its covariance decay is an equality,
and its projected transfer operator is proved nonzero.  Thus the construction
is an exact invariant-mode interface, not a claim that $\SU$ is finite and not
a construction of four-dimensional continuum Yang--Mills.  All decisive
statements compile in Lean 4 against pinned parent, satellite, and Mathlib
revisions.
\end{abstract}

\paragraph{Artifact statement.}
The formal source is
\texttt{YangMills/OS/SU2HeatTransport.lean}; its two concrete dependencies are
the pinned satellite \texttt{lean-2d-yang-mills} and the parent module
\texttt{SU2TransportWitness.lean}.  Exact revisions, byte counts, line-ending
hashes, build commands, and runtime information are recorded in the release
manifest distributed with this paper.

\section{Problem and contribution}

For a compact gauge group, heat-kernel convolution is the natural transfer
mechanism behind the exact solution of two-dimensional Yang--Mills theory.
At the formal level, however, three distinctions matter.

First, an abstract transfer theorem is not yet an inhabited gauge-theoretic
endpoint.  Second, a finite Markov matrix carrying a suggestive parameter is
not by itself a theorem about the Haar heat operator.  Third, a numerical
certificate for a matrix identity is not a Lean witness inhabiting $\SU$.
The purpose of the present construction is to close precisely this interface.

The contribution is the following chain, checked without a project-local
axiom and without assuming any of its conclusions:
\begin{enumerate}[label=(\roman*)]
\item finite heat-kernel partial sums act diagonally on every character already
  present in the sum;
\item dominated convergence upgrades this to the actual infinite heat kernel;
\item the Casimir ordering $3\le n(n+2)$ for $n\ge1$ identifies $\chi_1$ as the
  sharp slowest non-vacuum mode;
\item evaluation at the identity proves $\chi_1\ne0$ exactly;
\item the eigenvalue $e^{-3t/4}$ is identified with the parameter of a
  separately proved two-holonomy Dobrushin witness; and
\item the combined theorem returns both the continuous eigenfunction identity
  and a positive projected-operator gap bound with a nonzero projected
  operator.
\end{enumerate}

The novelty is deliberately narrower than the exact area-law results already
available in the satellite.  We do not rename those results as new.  Instead,
we add the missing spectral map needed by the parent transport endpoint.

\section{Canonical endpoint map}

Table~\ref{tab:endpoint} records the exact division of labor.  This is also the
minimum interpretation that connects the continuous heat kernel to the
published parent theorem: omitting the infinite convolution theorem leaves the
finite rate semantically ungrounded, while demanding a completed full
$L^2(\SU)$ spectral theorem would exceed what the endpoint consumes.

\begin{table}[ht]
\centering
\small
\begin{tabular}{p{.22\textwidth}p{.34\textwidth}p{.34\textwidth}}
\toprule
Layer & Exact object & Role \\
\midrule
Satellite & normalized Haar probability, Chebyshev characters, finite and
infinite class heat kernels, character convolution orthogonality & published
analytic input \\
New bridge & \texttt{heatKernel\_character\_eigen}, Casimir-rate ordering,
nonzero fundamental mode & continuous spectral interface \\
Frozen parent & two bundled $\SU$ holonomies, exact two-state covariance,
\texttt{exact\_transport\_witness} & finite non-vacuous endpoint \\
Abstract endpoint & abstract uniform-gap theorem & converts exact
band decay into a projected transfer norm bound \\
\bottomrule
\end{tabular}
\caption{Canonical endpoint map.  The middle row is the new theorem lane.}
\label{tab:endpoint}
\end{table}

The parent endpoint is not merely existential at the carrier level.  Its type
is a sigma type whose elements contain a finite label, a bundled $2\times2$
complex matrix, a proof of unitarity, a proof of determinant one, and an
equality with one of two explicitly constructed holonomies.  Hence every
carrier element inhabits $\SU$ in Lean.

\section{The concrete heat operator}

Let $\Haar$ be normalized Haar probability on $\SU$.  For functions $f,h$ the
satellite convolution convention is
\[
  (f*h)(g)=\int_{\SU} f(x)h(x^{-1}g)\,d\Haar(x).
\]
Write $\chi_n$ for the real-valued Chebyshev character of highest-weight label
$n$, regarded as a complex function, and set
\[
  a_n(t)=\exp\!\left(-\frac{t\,n(n+2)}4\right).
\]
The class heat kernel is the convergent series
\[
  K_t(g)=\sum_{n=0}^{\infty}(n+1)a_n(t)\chi_n(g).
\]
The factor $n+1$ is the representation dimension, while $n(n+2)/4$ is the
$\SU$ quadratic Casimir in the satellite's normalization.

The input orthogonality statement is already stronger than a scalar inner
product identity.  It gives the convolution formula
\begin{equation}\label{eq:charconv}
  (\chi_n*\chi_m)(g)=
  \begin{cases}
    \chi_n(g)/(n+1),&n=m,\\
    0,&n\ne m.
  \end{cases}
\end{equation}
This normalization cancels the dimension in $K_t$ exactly.

\subsection{Finite partial sums}

Let $K_t^{<N}=\sum_{n<N}(n+1)a_n(t)\chi_n$.  The first new Lean theorem is the
finite identity below.

\begin{theorem}[finite character eigenidentity]\label{thm:finite}
For $m<N$, $t\in\mathbb R$, and $g\in\SU$,
\[
  (K_t^{<N}*\chi_m)(g)=a_m(t)\chi_m(g).
\]
\end{theorem}

\begin{proof}
Finite integrability follows from continuity on the compact group.  Linearity
of the Haar integral turns convolution with the partial sum into a finite sum
of integrals.  Equation~\eqref{eq:charconv} kills every term except $n=m$;
the remaining factor $(m+1)$ cancels its reciprocal.  The Lean proof expands
the integrand before invoking the satellite theorem, so no desired
eigenidentity is passed as a premise.
\end{proof}

The formal statement is:
\begin{lstlisting}
theorem heatKernelPartial_character_eigen_of_lt
    (N : Nat) (t : Real) (m : Nat) (hm : m < N) (g : SU2) :
  su2Convolution (su2HeatKernelPartial N t)
      (su2CharacterChebyshev m) g =
    su2ClassHeatWeight t m * su2CharacterChebyshev m g
\end{lstlisting}

\subsection{Passage to the infinite kernel}

The infinite theorem requires more than rewriting a formal series.  For fixed
$m$ and $g$, define
\[
  F_N(x)=K_t^{<N}(x)\chi_m(x^{-1}g),\qquad
  F(x)=K_t(x)\chi_m(x^{-1}g).
\]
The satellite supplies pointwise convergence $K_t^{<N}(x)\to K_t(x)$ for
$t>0$ and a summable heat-kernel majorant.  It also supplies
$|\chi_m(y)|\le m+1$.  Thus
\[
  |F_N(x)|\le
  \left(\sum_{n=0}^{\infty}M_t(n)\right)(m+1),
\]
a constant integrable against the probability measure $\Haar$.  Dominated
convergence passes the limit through the Haar integral.  For all $N\ge m+1$,
Theorem~\ref{thm:finite} makes that integral the same constant; uniqueness of
limits completes the proof.

\begin{theorem}[infinite heat-mode identity]\label{thm:infinite}
If $t>0$, then for every $m\in\mathbb N$ and $g\in\SU$,
\[
  (K_t*\chi_m)(g)=
  \exp\!\left(-\frac{t\,m(m+2)}4\right)\chi_m(g).
\]
\end{theorem}

The crucial point is quantification over the actual bundled compact group and
the actual infinite kernel.  No finite carrier appears in
Theorem~\ref{thm:infinite}.

\section{Sharp fundamental mode}

Define $r_n(t)=\exp(-t n(n+2)/4)$.  The vacuum mode is $r_0(t)=1$.  For
$n\ge1$ the elementary inequality
\[
  3\le n(n+2)
\]
and monotonicity of the real exponential give
\[
  r_n(t)\le r_1(t)=e^{-3t/4}\qquad(t\ge0).
\]
For $t>0$, $0<r_1(t)<1$.  Consequently the fundamental representation is the
sharp slowest mode among all non-vacuum character labels.  ``Sharp'' here
means that equality is attained at $n=1$, not that a Peter--Weyl completeness
theorem for the completed class-function Hilbert space has been formalized.

Non-vacuity is exact.  The satellite proves $\chi_n(1)=n+1$; hence
$\chi_1(1)=2$, and function extensionality shows $\chi_1$ cannot be the zero
function.  The specialization of Theorem~\ref{thm:infinite} is therefore a
genuine eigenmode:
\begin{equation}\label{eq:fundamental}
  (K_t*\chi_1)(g)=e^{-3t/4}\chi_1(g),\qquad \chi_1\ne0.
\end{equation}

\section{Finite \texorpdfstring{$\SU$}{SU(2)}-inhabited transport witness}

The finite endpoint is a quotient/interface object, not an approximation in
which group axioms are discarded.  Its two labels carry two explicit matrices:
the identity and a noncentral diagonal phase with determinant one.  Both are
bundled as elements of $\SU$, and their inequality is proved exactly.

For a rate $0<r<1$, the transition kernel is
\[
  W_r=\frac12
  \begin{pmatrix}1+r&1-r\\1-r&1+r\end{pmatrix}.
\]
Its rows sum to one, it is symmetric and strictly positive, and
\[
  W_r^k=W_{r^k}.
\]
For any observable $f$ on the two carrier points, the band covariance is the
identity
\begin{equation}\label{eq:cov}
  \operatorname{Cov}_k(f,f)=
  \left(\frac{f(0)-f(1)}2\right)^2r^k.
\end{equation}
Thus the decay premise consumed by the Dobrushin theorem is derived, not
assumed.  The sign observable $(1,-1)$ is orthogonal to the normalized constant
vacuum, and the projected transfer sends its first component to $r$.  Since
$r>0$, the projected operator is not zero.

The frozen witness instantiates $r$ by
\[
  \operatorname{heatRate}(t)=e^{-3t/4}.
\]
The new bridge proves this is definitionally the same real number as $r_1(t)$
in the continuous heat theorem.  It then combines
Equation~\eqref{eq:fundamental} with the existing endpoint:

\begin{theorem}[continuous-to-finite transport bridge]\label{thm:bridge}
For $t>0$:
\begin{enumerate}[label=(\alph*)]
\item $K_t*\chi_1=\operatorname{heatRate}(t)\chi_1$ pointwise on $\SU$;
\item $\chi_1\ne0$;
\item there exists $m>0$ for which the norm of the vacuum-projected finite
transfer operator is at most $e^{-m}$; and
\item that projected operator is nonzero.
\end{enumerate}
\end{theorem}

The theorem does not claim that the finite matrix equals the full heat operator
on $L^2(\SU)$.  It states the smaller fact actually required by the endpoint:
the non-vacuum rate used in the exact finite quotient is the eigenvalue of a
nonzero mode of the genuine Haar heat operator.

\section{Formal proof architecture}

The formal dependency graph is intentionally short:
\[
\begin{array}{c}
\text{satellite character convolution + heat majorant}\\
\Downarrow\\
\texttt{heatKernelPartial\_character\_eigen\_of\_lt}\\
\Downarrow\ \text{dominated convergence}\\
\texttt{heatKernel\_character\_eigen}\\
\Downarrow\ \text{Casimir ordering + }\chi_1(1)=2\\
\texttt{heatKernel\_fundamental\_eigen}\\
\Downarrow\ \text{rate identification}\\
\texttt{continuous\_mode\_to\_finite\_transport}.
\end{array}
\]

Table~\ref{tab:hypotheses} performs the hypothesis-discharge audit requested
for the endpoint.

\begin{table}[ht]
\centering
\small
\begin{tabular}{p{.25\textwidth}p{.28\textwidth}p{.36\textwidth}}
\toprule
Required fact & Source & Discharge \\
\midrule
$t>0$ & bridge input & used for summability and $r<1$ \\
Haar integrability & compact continuity & satellite continuity lemmas and
probability measure \\
dominating function & heat majorant & summable constant times $m+1$ \\
character orthogonality & satellite theorem & used termwise before the limit \\
Casimir ordering & new arithmetic lemma & proved by casting $1\le n$ and
nonlinear arithmetic \\
$\chi_1\ne0$ & new exact lemma & evaluation at the identity gives $2$ \\
finite kernel laws & frozen witness & explicit two-by-two algebra \\
decay premise & frozen witness & equality~\eqref{eq:cov}, not a hypothesis \\
projected nonzero & frozen witness & exact sign-vector evaluation \\
\bottomrule
\end{tabular}
\caption{Every load-bearing hypothesis and where it is discharged.}
\label{tab:hypotheses}
\end{table}

The Lean proof uses classical choice and standard analytic facts inherited
from Mathlib, as expected for Haar integration and infinite series.  The audit
criterion is therefore not an empty \texttt{\#print axioms} list, but absence
of \texttt{sorryAx} and absence of project-local axioms.  The exact printed
axiom sets are captured in the verification log.

\section{Reproducibility and frozen revisions}

The artifact uses Lean 4 through the parent repository's pinned toolchain.
The satellite is fetched as an exact Lake dependency, not as an unversioned
copy.  At fabrication time the pins were:
\begin{center}
\begin{tabular}{ll}
\toprule
Object & Revision \\
\midrule
parent starting point & \texttt{e6b672dda3e41da25a22178a7ef73ab8e7f8edb7} \\
\texttt{lean-2d-yang-mills} & \texttt{05c4ec316cb9aa295416670a2578b1c2e77e1c36} \\
Mathlib & \texttt{07642720480157414db592fa85b626dafb71355b} \\
\bottomrule
\end{tabular}
\end{center}

The decisive commands are:
\begin{lstlisting}[language={}]
lake update
lake build YangMills.OS.SU2HeatTransport
lake build YangMillsCore
lake env lean Task48Axioms.lean
\end{lstlisting}
They were run in a fresh Google Colab Pro+ CPU/high-RAM runtime with no GPU.
Local desktop work was limited to source editing, Git, hashing, TeX compilation
and PDF rendering.  A source file on disk was not counted as verified until its
\texttt{.olean} had been materialized remotely.

\section{Scope, limitations, and interpretation}

The results support the following exact statement: the infinite class heat
kernel has a nonzero fundamental eigencharacter at rate $e^{-3t/4}$, and this
same rate parametrizes a genuine $\SU$-inhabited finite transfer witness that
fires the parent gap endpoint non-vacuously.

They do not establish any of the following:
\begin{itemize}
\item four-dimensional continuum Yang--Mills or the Clay mass-gap problem;
\item a thermodynamic or continuum limit of the finite witness;
\item Peter--Weyl completeness or an operator-norm computation on the completed
full Haar $L^2$ class space;
\item a new proof of the satellite's Migdal recursion, semigroup law, or exact
finite-cellulation area law; or
\item that two carrier points exhaust, discretize, or approximate the compact
group.
\end{itemize}

The right interpretation is an exact transport interface.  It replaces a
previously merely rate-labelled finite endpoint by a theorem that its rate is
realized by a nonzero mode of the continuous operator.  This is sufficient for
the published parent theorem and no stronger completion claim is needed.

\section{Relation to the literature}

The use of heat kernels and character expansions in two-dimensional
Yang--Mills goes back to Migdal's recursion equations~\cite{Migdal1975}.
Driver constructed the Yang--Mills measure and related it to the heat equation
on compact groups~\cite{Driver1989}; rigorous compact-surface developments
include Sengupta~\cite{Sengupta1997} and L\'evy~\cite{Levy2003}.  The present
paper does not reprove those broad analytic constructions.  It contributes a
small, exact, inspectable theorem interface inside Lean.

The formalization builds on Mathlib~\cite{Mathlib2020} and on the pinned
two-dimensional satellite~\cite{Lean2dYM}.  The repository citation is a
commit permalink: it identifies the precise source used by Lake rather than a
moving branch.

\section{Conclusion}

The endpoint is now connected at the correct level of abstraction.  The
continuous side is an actual Haar integral on bundled $\SU$, the infinite
series is handled by dominated convergence, the fundamental eigencharacter is
proved nonzero, and the finite side consists of distinct genuine holonomies
with an exact covariance identity and nonzero projected transport.  The rate
$e^{-3t/4}$ is no longer an analogy between two modules; it is a checked
equality used in a single end-to-end Lean theorem.

\begin{thebibliography}{9}

\bibitem{Migdal1975}
A. A. Migdal,
``Recursion equations in gauge field theories,''
\emph{Soviet Physics JETP} 42 (1975), 413--418.
\url{http://www.jetp.ras.ru/cgi-bin/e/index/e/42/3/p413?a=list}

\bibitem{Driver1989}
B. K. Driver,
``YM$_2$: continuum expectations, lattice convergence, and lassos,''
\emph{Communications in Mathematical Physics} 123 (1989), 575--616.
\href{https://doi.org/10.1007/BF01238812}{doi:10.1007/BF01238812}.

\bibitem{Sengupta1997}
A. Sengupta,
\emph{Gauge Theory on Compact Surfaces},
Memoirs of the American Mathematical Society 126, no. 600 (1997).
\href{https://doi.org/10.1090/memo/0600}{doi:10.1090/memo/0600}.

\bibitem{Levy2003}
T. L\'evy,
\emph{Yang--Mills Measure on Compact Surfaces},
Memoirs of the American Mathematical Society 166, no. 790 (2003).
\href{https://doi.org/10.1090/memo/0790}{doi:10.1090/memo/0790}.

\bibitem{Mathlib2020}
The mathlib Community,
``The Lean Mathematical Library,''
in \emph{Proceedings of CPP 2020}, 367--381.
\href{https://doi.org/10.1145/3372885.3373824}{doi:10.1145/3372885.3373824}.

\bibitem{Lean2dYM}
L. Eriksson,
\emph{lean-2d-yang-mills}, commit
\href{https://github.com/lluiseriksson/lean-2d-yang-mills/tree/05c4ec316cb9aa295416670a2578b1c2e77e1c36}
{\texttt{05c4ec316cb9aa295416670a2578b1c2e77e1c36}} (2026).

\end{thebibliography}

\end{document}
